\subsubsection*{S5: generated layouts}
Across the crossed drone-count/noise set, $\val{RandOverall}$ episodes
succeed. At $\alpha\le0.25$, $\val{RandLowNoise}$ runs succeed,
including $\val{RandNFourLow}$ with four drones; four drones at $\alpha=2$
succeed in $\val{RandNFourHigh}$. This indicates sensitivity to noise in
the tested layouts, while increasing $N$ jointly enlarges state and action
spaces. Mean new-node depth is $\val{DepthN1}$ for one drone and
$\val{DepthN4}$ for four, consistent with less successor reuse when joint
outcomes multiply. Each $(N,\alpha)$ pair has only one layout, so the
per-layout results in Appendix~\ref{app:layouts} do not isolate layout effects.
For the four-drone layout at $\alpha=2$, joint value iteration was not attempted
($140^4\approx4\times10^8$ joint positions per step). Each drone alone could
achieve at least $\val{CeilFourSingle}\%$ success when other drones are removed.
Executing those policies together gives $\val{CeilFourUncoord}\%$ success
(1,000 diagnostic runs; Wilson 95\% interval $\val{CeilFourUncoordCI}$),
compared with $\val{RandNFourHigh}$ for the solver. The different sample sizes
and unpaired comparison do not establish a solver advantage or isolate the
cause of failure. They motivate further coordination tests
(\texttt{results/ceilings.json}).
